\usepackage{amssymb}
\usepackage{amsmath}
\usepackage{times}
\usepackage{psfrag}
\usepackage{amssymb}
\usepackage{wasysym}
\usepackage{float}
%\usepackage{esz}
\usepackage{listings}

\usepackage{afterpage}

\usepackage{graphicx}
\usepackage{multirow}
\restylefloat{table}
\usepackage{listings}
\newcommand{\lstocl}{\lstset{emph={set,bag,or,and,not,implies,select,includes,let,in,forAll,context,inv,self,pre,post,if,then,else,endif},emphstyle=\textbf,basicstyle=\scriptsize,commentstyle=\scriptsize,showstringspaces=false,frame=single,breaklines=true,breakatwhitespace=true,frame=tlrb,numbers=none,framexleftmargin=0mm,numbersep=2mm}}

\setlength{\parindent}{0mm}
\usepackage{fancybox} 
\usepackage{fancyheadings} %Header
\usepackage{amsmath}
%\usepackage{amsmath,amssymb,amstext,stmaryrd} %Zusätzlichezeichen
%\usepackage{float}
\usepackage{graphicx}%Bilder
\usepackage{ifthen}%Solution switch
\usepackage[ngerman]{babel}%Sprachpacket für Silbentrennung
\usepackage[utf8]{inputenc}%Sonderzeichen
%\usepackage[T1]{fontenc}

% \usepackage{verbatim}
% \usepackage{moreverb}
\usepackage{color}%Farben
\usepackage[
        pdftitle={Exam of the course ``\course'' \semester{}},
        pdfsubject={Exam},
        pdfpagemode=None,
        a4paper=true,
        colorlinks=false,
        urlcolor=blue]{hyperref}%PDF titel
%\usepackage{oz}%Z und OZ Formatierung
%\usepackage{tikz}%Hübsche Bilder
%\usetikzlibrary{decorations.pathmorphing,shapes,decorations.text,decorations,calc,positioning,automata,matrix,arrows,chains}

\usepackage[normalem]{ulem}

\newcommand{\taskdescription}[2]{#1 \\ \emph{#2}}


%% Einrückung der Aufzählungen verkleinern
%\setlength\leftmargini{0pt}
%\setlength\leftmargin{5ex}
\setlength\leftmarginii{20pt}
\setlength\leftmargini{0pt}

%% Kommandas für die Lösungen
\newcommand{\sol}[2]{\ifthenelse{\equal{\solution}{true}}{\textcolor{red}{\newline Solution:\\#1}}{#2}}
\newcommand{\solOCL}[2]{\ifthenelse{\equal{\solution}{true}}{\newline Solution:#1}{#2}}
\newcommand{\solTable}[2]{\ifthenelse{\equal{\solution}{true}}{\textcolor{red}{#1}}{#2}}
\newcommand{\solText}[2]{\ifthenelse{\equal{\solution}{true}}{\textcolor{red}{#1}}{#2}}

%% Es dürfen keine Leerzeichen vor einem \begin{solu} oder \end{solu} sein
\usepackage{comment}
\specialcomment{solu}{\begingroup~\\\color{red}\itshape}{\endgroup}
\newenvironment{nonsolu}{}{}
\ifthenelse{\equal{\solution}{false}}{\excludecomment{solu}}{\excludecomment{nonsolu}}

\newcommand{\solRule}[2]{\ifthenelse{\equal{\solution}{false}}
{\rule{#2}{0.1mm}}{\textcolor{red}{#1}}}

\renewcommand{\labelenumi}{(\arabic{enumi})}
\renewcommand{\labelenumii}{(\arabic{enumi}.\arabic{enumii})}

%Kopf- und Fußzeilen setzen
\pagestyle{fancy}
%\setlength{\headrulewidth}{0.0pt}
%\setlength{\footrulewidth}{0.4pt}

\newcommand{\centeredgraphic}[1]{\centerline{\resizebox{0.8\linewidth}{!}{\includegraphics[angle=270]{#1}}}}

% Punktezähler
\newcounter{overallpoints}%[counter] 
\newcommand{\points}[1]{
\ifthenelse{\equal{#1}{1}}
{~\\ \marginpar{ / #1 pt} }
{~\\ \marginpar{ / #1 pts}  }
\addtocounter{overallpoints}{#1} }

\ifthenelse{\equal{\solution}{false}}
{
\newcommand{\solListing}{f}
}{
\newcommand{\solListing}{t}
}


\newcommand{\squareM}{
\ifthenelse{\equal{\solution}{false}}
{\square}
{ \textcolor{red}{\boxtimes} } 
}

\ifthenelse{\equal{\solution}{false}}
{
\newcommand{\emphc}[1]{#1}
\newcommand{\textbfc}[1]{#1}
}{
\newcommand{\emphc}[1]{\emph{#1}}
\newcommand{\textbfc}[1]{\textbf{#1}}
}

% \newcommand{\sol}[2]{\ifthenelse{\equal{\solution}{false}}
% {\rule{#2}{0.1mm}}{\textcolor{red}{#1}}}

\newcommand{\sollineOne}[1]{\newline\ifthenelse{\equal{\solution}{false}}
{\rule{\linewidth}{0.1mm}}{\textcolor{red}{#1}}}

\newcommand{\sollineTwo}[1]{\newline ~\\ \ifthenelse{\equal{\solution}{false}}
{\rule{\linewidth}{0.1mm} \\[0.5ex]
\rule{\linewidth}{0.1mm}}{\textcolor{red}{#1}}}

\newcommand{\pointassign}[1]{
\ifthenelse{\equal{\solution}{true}}
{ \textcolor{red}{#1} } 
{}
}

\lstset{language=Java,basicstyle=\footnotesize,showstringspaces=false} %basicstyle=\small

\newcommand{\wrong}[1]{}%{\ifthenelse{\equal{\solution}{false}}{}{\textcolor{blue}{\sout{#1}}}}

\usepackage[sf,bf]{titlesec}
\titleformat{\section}[hang]{\large\fontseries{bx}\selectfont}{\thesection}{5mm}{}[\vspace{0mm}]
\titleformat{\subsection}[hang]{\normalsize\fontseries{bx}\selectfont}{\thesubsection}{5mm}{}[\vspace{0mm}]
\setlength{\textheight}{240mm} 

\usepackage{fancyhdr}
\pagestyle{fancy}

\lhead{}
\chead{}
\rhead{}
\cfoot{}
%\lfoot{\textit{Name:} \hfill \textit{Student ID:}\hspace*{50mm}}
%\rfoot{Page - \thepage\ -}
\renewcommand{\headrulewidth}{0.0pt}
\renewcommand{\footrulewidth}{0.4pt}

